\documentclass[11pt,a4paper]{article}
\usepackage{notes}

\title{Solution Set: Gravity Anomalies and Isostatic Response}
\author{}
\date{}

\begin{document}
\maketitle

%%%%%%%%%%%%%%%%%%%%%%%%%%%%%%%%%%%%%%%%%%%%%%%%%%%%%%%%%%%%
\section*{Question 2: Shallow and Deep Bodies}

\subsection*{(A) Gravity anomaly of a 2-D polygon}

The Talwani et al.\ (1959) formulation gives the exact vertical gravity anomaly for a
two-dimensional body of constant density contrast $\Delta\rho$ with polygonal geometry.
The method assumes:
\begin{itemize}
    \item infinite extent perpendicular to the $x$--$z$ plane,
    \item constant density contrast,
    \item observation at the surface ($z=0$).
\end{itemize}

The total anomaly is obtained as the sum of contributions from each polygon edge.
No approximation is made beyond the 2-D assumption. Implementation is most easily
performed by looping over polygon edges in Python or MATLAB.

%%%%%%%%%%%%%%%%%%%%%%%%%%%%%%%%%%%%%%%%%%%%%%%%%%%%%%%%%%%%
\subsection*{(B) Triangular rift basin}

\subsubsection*{(1) Geometry}

The basin is bounded by two faults dipping at $60^\circ$ and has a total vertical
offset of $\SI{5}{km}$. The basin is centered at $(0,0)$, with the surface at $z=0$.

The horizontal half-width of the basin is:
\[
w = \frac{(5/2)}{\tan 60^\circ} \approx \SI{1.44}{km}.
\]

A suitable clockwise ordering of vertices is:
\[
(-1.44,0), \quad (0,-5), \quad (1.44,0).
\]

\subsubsection*{(2) Gravity anomaly}

Using the polygon equation with a negative density contrast typical of sediments
($\Delta\rho \approx -400~\si{kg\,m^{-3}}$), the computed anomaly is negative and symmetric.
The anomaly decays smoothly to zero at large distances from the basin.

\subsubsection*{(3) Interpretation}

\begin{itemize}
    \item The maximum horizontal gradient occurs near the basin edges, not at the basin center.
    \item Because the geometry is symmetric, the anomaly is symmetric about $x=0$.
\end{itemize}

The gravity signal is dominated by the sloping fault planes rather than the deepest
point of the basin.

%%%%%%%%%%%%%%%%%%%%%%%%%%%%%%%%%%%%%%%%%%%%%%%%%%%%%%%%%%%%
\subsection*{(C) Trapezoidal rift basin}

\subsubsection*{(1) Geometry}

The trapezoidal basin has the same surface coordinates as part (B) but a flat base
at $\SI{3}{km}$ depth. The sidewall dips are comparable to those in the triangular case.
One valid vertex ordering is:
\[
(-1.44,0),\ (-0.44,-3),\ (0.44,-3),\ (1.44,0).
\]

\subsubsection*{(2) Gravity anomaly}

The trapezoidal basin produces a broader, smoother anomaly compared to the triangular
basin, with reduced curvature near the edges.

\subsubsection*{(3) Comparison}

\begin{itemize}
    \item \textbf{Edge slope} controls the location and sharpness of gradients.
    \item \textbf{Depth extent} controls anomaly amplitude and wavelength.
    \item \textbf{Density contrast} scales the magnitude but does not strongly affect shape.
\end{itemize}

Different basin geometries with similar edges and depths can produce very similar
gravity anomalies.

%%%%%%%%%%%%%%%%%%%%%%%%%%%%%%%%%%%%%%%%%%%%%%%%%%%%%%%%%%%%
\subsection*{(D) Basin over a crustal root}

The crustal root is represented by a triangular body of width $\SI{30}{km}$ and
maximum thickness $\SI{10}{km}$. Typical density contrasts are:
\[
\Delta\rho_{\text{basin}} \sim -400~\si{kg\,m^{-3}}, \qquad
\Delta\rho_{\text{root}} \sim +300~\si{kg\,m^{-3}}.
\]

\subsubsection*{Root alone}

The root produces a long-wavelength, low-amplitude anomaly dominated by its depth
and width.

\subsubsection*{Basin alone}

The basin produces a short-wavelength anomaly dominated by edge effects.

\subsubsection*{Combined model}

When the basin and root are combined, the shallow basin dominates the short-wavelength
signal, while the root affects the long-wavelength background. With suitable basin
parameters, the combined anomaly can closely resemble a basin-only anomaly.

This demonstrates the fundamental non-uniqueness of gravity interpretation.

%%%%%%%%%%%%%%%%%%%%%%%%%%%%%%%%%%%%%%%%%%%%%%%%%%%%%%%%%%%%
\section*{Question 3: Isostasy --- Oceanic vs Continental Lithosphere}

Isostatic columns constructed using crustal thickness alone cannot balance oceanic
and continental lithosphere. Balance requires a density difference within the mantle
lithosphere of order:
\[
\Delta\rho \sim 30\text{--}60~\si{kg\,m^{-3}},
\]
which can plausibly be explained by compositional or thermal differences.

%%%%%%%%%%%%%%%%%%%%%%%%%%%%%%%%%%%%%%%%%%%%%%%%%%%%%%%%%%%%
\section*{Question 4: Isostatic Response to Erosion}

\subsection*{(A) Single-layer crust}

Assume:
\[
\rho_c = 2800~\si{kg\,m^{-3}}, \qquad \rho_m = 3300~\si{kg\,m^{-3}}.
\]

The Airy isostatic factor is:
\[
\alpha = \frac{\rho_m - \rho_c}{\rho_m} \approx 0.15.
\]

\subsubsection*{(1) Instantaneous erosion of \SI{2}{km}}

Surface lowering:
\[
\Delta \varepsilon = 2~\si{km} \times \alpha \approx 0.30~\si{km}.
\]
New elevation:
\[
\varepsilon \approx \SI{5.8}{km}.
\]

\subsubsection*{(2) Total erosion to baseline}

\[
E = \frac{\varepsilon_0}{\alpha} \approx \frac{6.1}{0.15} \approx \SI{41}{km}.
\]

\subsubsection*{(3) Fraction replaced by rebound}

\[
1-\alpha \approx 85\%.
\]

%%%%%%%%%%%%%%%%%%%%%%%%%%%%%%%%%%%%%%%%%%%%%%%%%%%%%%%%%%%%
\subsection*{(B) Continuous erosion with rebound}

The erosion law is:
\[
\frac{dE}{dt} = k\,\varepsilon,
\qquad k = \SI{0.1}{Myr^{-1}}.
\]

Under Airy isostasy:
\[
\frac{d\varepsilon}{dt} = -\alpha \frac{dE}{dt}
= -\alpha k \varepsilon.
\]

\subsubsection*{(1) General solution}

\[
\varepsilon(t) = \varepsilon_0 e^{-\alpha k t}.
\]

\subsubsection*{(2) No rebound limit}

If $\alpha=1$:
\[
\varepsilon(t) = \varepsilon_0 e^{-k t}.
\]

\subsubsection*{(3) Initial erosion rates}

With $\varepsilon_0 = \SI{6100}{m}$:
\[
\left.\frac{dE}{dt}\right|_{t=0} = 610~\si{m\,Myr^{-1}}, \qquad
\left.\frac{d\varepsilon}{dt}\right|_{t=0} \approx 91~\si{m\,Myr^{-1}}.
\]

\subsubsection*{(4) Time to erode 95\% of relief}

\[
t_{95} = \frac{\ln(20)}{\alpha k}
\approx \SI{200}{Myr}.
\]

\subsubsection*{(5) Interpretation}

Isostatic rebound increases erosion timescales by more than an order of magnitude,
allowing large mountain belts to persist for hundreds of millions of years despite
vigorous erosion.

\end{document}